
\documentclass[11pt]{article}

\usepackage{amsfonts}
\usepackage{amsmath}
\usepackage{amscd}
\usepackage{amssymb}
\usepackage[english]{babel}
\usepackage[UTF8]{ctex}
\usepackage[tableposition=top]{caption}
\usepackage{indentfirst}
\usepackage{natbib}
\usepackage{rotating}
\usepackage[colorlinks=true,allcolors=blue]{hyperref}
\usepackage{url}
\usepackage{doi}
\usepackage{tikz-cd}

\DeclareMathOperator{\idfun}{id}
\DeclareMathOperator{\isContr}{isContr}
\DeclareMathOperator{\isProp}{isProp}
\DeclareMathOperator{\isSet}{isSet}
\DeclareMathOperator{\dom}{dom}
\DeclareMathOperator{\cod}{cod}
\DeclareMathOperator{\qinv}{qinv}
\DeclareMathOperator{\binv}{biinv}
\DeclareMathOperator{\isEquiv}{isequiv}
\DeclareMathOperator{\refl}{refl}
\DeclareMathOperator{\inl}{inl}
\DeclareMathOperator{\inr}{inr}
\DeclareMathOperator{\successor}{succ}
\DeclareMathOperator{\inductor}{ind}
\DeclareMathOperator{\recursor}{rec}

\newcommand{\sloop}{\text{loop}}
\newcommand{\sbase}{\text{base}}
\newcommand{\ssurf}{\text{surf}}

\newcommand{\set}[1]{\{\,#1\,\}}
\newcommand{\proptrunc}[1]{\lVert #1 \rVert}
\newcommand{\termtrunc}[1]{\lvert #1 \rvert}
\newcommand{\bigproptrunc}[1]{\left\lVert #1 \right\rVert}

\newcommand{\ints}{\mathbb{Z}}
\newcommand{\nats}{\mathbb{N}}
\newcommand{\real}{\mathbb{R}}
\newcommand{\universe}{\mathcal{U}}
\newcommand{\rats}{\mathbb{Q}}

\newcommand{\fatdot}{\,\cdot\,}

\let\emptyset=\varnothing

\newcommand{\REVISED}{\begin{center} \LARGE REVISED DOWN TO HERE \end{center}}
\newcommand{\MOVED}[1][equation]{\begin{center} [#1 moved] \end{center}}

\begin{document}

\title{Reply to Ladyman and Presnell (2018)}
\author{Charles J. Geyer}
\maketitle

\section{License}

This work is licensed under a Creative Commons
Attribution-ShareAlike 4.0 International License
(\url{http://creativecommons.org/licenses/by-sa/4.0/}).

The \LaTeX\ source for this document is
at \url{https://github.com/cjgeyer/Survey}.

\section{Introduction}

This document is a comment on \citet{ladyman-presnell-foundations}.

\section{Criteria for Foundation}

\citet{ladyman-presnell-foundations} present five criteria for a framework
such as homotopy type theory (HoTT) to be a proper philosophical foundation
for mathematics.
\begin{enumerate}
\item A single framework in which to cast some or all of existing mathematics.
\item A semantics saying how theoretical terms are understood and how rules
    of inference are understood.
\item A metaphysics that spells out the ontological status of entities.
\item An epistemology.
\item Methodology based on the above.
\end{enumerate}

Item~1 is very weak.  Some does not require much.  They also insist that
the foundation can be very impractical.  My paper argues that first order
classical logic plus ZFC (Zermelo–Fraenkel set theory with the axiom of
choice) is (with 120 years of hindsight) a horrible foundation for
mathematics.  That does not bother them even a little bit.

Item~3 is important to some philosophers but not to mathematicians,
who consider it BS they don't need.  I personally say that metaphysics
is a disease given to speakers of Indo-European languages with their
curious fetishization of the verb to be.  Indo-European language speakers
are always speaking about what is is.  Other humans not so much.

HoTT is concerned with how math works, what you can do with it.

Item~4 is blatently obvious for HoTT.  You can verify with proof assistants
(Rocq, Agda, etc.).

Item~5 I really agree with.  What we want is a \emph{dao} (道) of mathematics.
Voevodsky wanted the same thing and had gone a long way towards developing
HoTT when he died.

Section~2.2 of \citet{ladyman-presnell-foundations} argues about autonomy.
If some purported foundation relies on some other foundation (is expressed
in terms of), then it is not foundational at all.
They note that some other authors (Linnebo and Pettigrew, 2011,
which they cite, but we do not)
assert that category theory is not foundational because
``standard presentations'' rely on set theory.  This is simple ignorance
of authoritative presentations.  Since \citet{lawvere-cat-of-cat} it has
been clear that category theory can be considered a foundation of mathematics
in no way dependent on set theory and perhaps in no way dependent on logic
\citep{lawvere-thesis}.  \citet{maclane} also makes this clear.

\citet{ladyman-presnell-foundations} say that HoTT is not autonomous,
at least as presented in 
\citet[hereinafter referred to as the HoTT book]{hott-book},
\defcitealias{hott-book}{HoTT book}
because it is dependent on homotopy theory, and classical homotopy theory
is dependent on large parts of mainstream mathematics, which is supposedly
founded on classical first-order logic plus ZFC.  But this is a mistake.
HoTT is Martin-L\"{o}f type theory (MLTT) \citep{martin-lof-mltt}
plus a few additions (the univalence axiom, homotopy levels, higher
inductive types), and there is no trace of homotopy in MLTT.
One might claim that homotopy levels and the univalence axiom are dependent
on homotopy, but
\begin{itemize}
\item if so, they are dependent on \emph{synthetic} homotopy theory (as
    reconstructed in HoTT) not on classical homotopy theory,
\item the univalence axiom can be considered to just be Leibniz's principle
    of the identity of indiscernibles (or its converse, if you consider that
    a different principle) applied to the universe (so not
    dependent on homotopy),
\item homotopy levels are there primarily to clarify the nature of mere
    propositions and mere sets in MLTT (so not dependent on homotopy),
\item homotopy is mostly mere analogy not part of the mathematics any more
    than diagrams in Euclidean geometry are part of the proofs, and
\item HoTT is for all of mathematics, not just homotopy.
\end{itemize}

Uffda!  \citet{ladyman-presnell-foundations} do not mention the univalence
axiom or homotopy levels or higher inductive types.  Thus they do not even
approach the difficult issues.
They are not even introducing the issues that make HoTT different from MLTT.
Hence their argument cannot even get off the ground.  MLTT is clearly not
in any way even influenced by homotopy.  This is just another case of
a piece of mathematics being far more useful than its creator imagined,
like Riemannian geometry being useful for general relativity or Hilbert spaces
for quantum mechanics.

\citet{ladyman-presnell-foundations} continue this point in their Section~4
where they continue to claim that HoTT is dependent on the homotopy
interpretation.  But there are three analogies to terms and types
\begin{table}
\caption{Analogies for Terms and Types}
\label{tab:term-type}
\begin{center}
\begin{tabular}{ll}
term & type \\
\hline
proof & proposition \\
element & set \\
point & space
\end{tabular}
\end{center}
\end{table}
shown in our Table~\ref{tab:term-type}.  \citet{ladyman-presnell-foundations}
are aware of all of these, spending all of their Section~3.3 on propositions
as types.  But the homotopy analogy (types as synthetic topological spaces,
terms as points in those spaces, and identifications as paths in those spaces)
is no more ``fundamental'' than the other two analogies.  There are just
terms and types (and identity types and their terms),
no matter how you woof about them.  The analogies are there
to help mathematicians familiar with conventional mathematics learn HoTT.
You could describe HoTT completely without mentioning them (as in the
appendicies of the \citetalias{hott-book}).  But that would just make it more
abstract and hence more mysterious to users.  The analogies are not
fundamentally part of HoTT.

\citet{ladyman-presnell-foundations} eventually (their Section~5.2) get
around to an analogy (interpretation) that suits them: types are general
concepts and terms are specific concepts.  This is not too bad, but we
can do a lot better by looking at the type systems of pure functional
programming languages like SML, OCaml, and Haskell.  A type is just a type.
And we know what types are.  They are kinds of objects that variables can
refer to.  And \emph{pace} \citet[Section~5.1]{ladyman-presnell-foundations}
this ``commits us to a strong kind of Platonism'' is wrong: this is just
another
analogy.  Indo-European language speakers like the verb ``to be.''  They like
to say ``this is that.''  And they are going to do that with mathematical
``objects'' (in scare quotes) whatever philosophy of mathematics they
may have, just like they do with ``objects'' (scare quotes again)
in computer programs.  Using the verb ``to be'' is not a philosophical
commitment.

\begin{thebibliography}{}

\bibitem[Ladyman and Presnell(2018)]{ladyman-presnell-foundations}
Ladyman, J., and Presnell, S. (2018).
\newblock Does homotopy type theory provide a foundation for mathematics?
\newblock \emph{The British Journal for the Philosophy of Science},
    \textbf{69}, 377--420.
\newblock \doi{10.1093/bjps/axw006}.

\bibitem[Lawvere(1963)]{lawvere-thesis}
Lawvere, F. W. (1963).
\newblock Functorial Semantics of Algebraic Theories.
\newblock Ph.D. thesis, Columbia University.
\newblock Republished (2004) with an author’s comment and a supplement in
    \emph{Reprints in Theory and Applications of Categories}, \textbf{5},
    1--121.
\newblock \url{http://www.tac.mta.ca/tac/reprints/articles/5/tr5.pdf}.

\bibitem[Lawvere(1966)]{lawvere-cat-of-cat}
Lawvere, F. W. (1966).
\newblock The category of categories as a foundation for mathematics.
\newblock In S. Eilenberg, D. K. Harrison, S. MacLane, H. Röhrl (eds.),
    \emph{Proceedings of the Conference on Categorical Algebra - La Jolla
    1965}, 1--20.
\newblock Springer.
\newblock \doi{10.1007/978-3-642-99902-4}.

\bibitem[Mac Lane(1998)]{maclane}
Mac Lane, S. (1998).
\newblock \emph{Categories for the Working Mathematician}, second edition.
\newblock Springer-Verlag, New York.

\bibitem[Martin-L\"{o}f(1975)]{martin-lof-mltt}
Martin-L\"{o}f, P. (1975).
\newblock An intuitionistic theory of types: predicative part.
\newblock In: H. E. Rose, J. C. Shepherdson (eds.),
    \emph{Logic Colloquium ‘73, Proceedings of the Logic Colloquium,
    Studies in Logic and the Foundations of Mathematics}, \textbf{80}, 73--118.
\newblock Elsevier.
\newblock \doi{10.1016/S0049-237X(08)71945-1}.

\bibitem[The Univalent Foundations Program(2013)]{hott-book}
The Univalent Foundations Program (2013).
\newblock \emph{Homotopy Type Theory: Univalent Foundations of Mathematics}.
\newblock Institute for Advanced Study, Princeton, NJ.
\newblock \url{https://homotopytypetheory.org/book}.

\end{thebibliography}

\end{document}

